\section{Limitations and Conclusion}
\label{sec:limitations}

\paragraph{Limitations.} As noted in Section~\ref{sec:setup}, absolute Recall in our from-scratch pipeline trails published GENIUS numbers; this sits on exactly the dataset and direction carrying our one clear positive result (MSCOCO Dir.\ A), so we cannot rule out a different effect size in a fully converged, published-parity regime -- a limitation we regard as consistent with, rather than contradicting, this paper's central claim that regularization's effect is not a stable, portable property of the method. Our significance tests certify robustness to decoder-training-seed noise, not to every source of variance: MSCOCO's gain is additionally verified across 3 independently retrained Stage-1 tokenizers, but VisualNews's reversal was checked only across decoder seeds on a single tokenizer per variant, so tokenizer-initialization variance is unmeasured there. Section~\ref{sec:finding3}'s correlations use four $\lambda$ points per dataset; the divergence in their behaviour between datasets is the finding, not $n{=}4$ significance. We localize the image$\to$text collapse to the tokenizer but do not fully explain it: text candidates start more anisotropic than image candidates in the fused embedding space (mean pairwise cosine similarity $0.842$ vs.\ $0.675$)~\cite{ethayarajh2019anisotropy}, which is consistent with per-sample balance pressure shattering an already-narrow text cluster first, but does not explain why the collapse persists at strong $\lambda$ once code entropy partially recovers -- a full mechanistic account is left to future work. Retrieval-tie-breaking within collision-heavy codes does not inflate MSCOCO's headline gain (verified against a random and an adversarial tie-break baseline, both within $0.05$ points of the reported value). Our three datasets are contrasting exemplars, not a sweep -- the optimal $\lambda$ for one dataset need not transfer to a further regime, which is itself part of what we report.

\paragraph{Conclusion.} Balance regularization is neither useless nor universally beneficial for RQ-based generative multimodal retrieval: it significantly helps MSCOCO text-to-image Recall while destroying image-to-text Recall on the same checkpoints, substantially hurts VisualNews's text-to-image Recall, and hurts FashionIQ outright. Decoupling regularization per target modality does not rescue either direction, because image and text candidates share one encoder and codebook. Both the direction-asymmetry and the dataset-dependence are already present at the RQ tokenizer level, before any decoder is trained. Practitioners should not adopt balance regularization by default without validating its effect per retrieval direction and per dataset.

\paragraph{Code and data.} Code, configurations, and per-seed results for every variant reported here will be released as a public repository at acceptance.
